\section{近似 BPO 更新场、梯度误差与势函数边界}
\label{opd:approximate-field}

生产路径先从精确平方回归推导 score 梯度，再以粗粒化中心替换数值权重。每步 stopped surrogate 的反向传播可计算该步向量；这一构造不保证所有参数位置的向量都是某个固定标量目标的梯度。

\begin{theorem}[粗粒化中心替换后的更新场]
\label{opd:field-curl}
固定 $q,t$ 和头部 $H$，记 $K=\KL(q\Vert p)$、$\epsilon_S=K-\widehat K_S$、$\epsilon_T=K_{qT}-\widehat K_T$。生产 BPO 的条件期望 score 更新为
\begin{equation}
\widehat g=\nabla L_{\mathrm{BPO}}-\epsilon_T\nabla K+\frac{\epsilon_S}{\eta}\nabla K.
\end{equation}
以 Jacobian 约定 $Dg_{ij}=\partial g_i/\partial\theta_j$，其反对称部分为
\begin{equation}
D\widehat g-(D\widehat g)^\top
=\eta^{-1}\left[\nabla K(\nabla\epsilon_S)^\top-\nabla\epsilon_S(\nabla K)^\top\right].
\end{equation}
该矩阵非零时，邻域内不存在二次连续可微的固定标量函数，其欧氏梯度等于 $\widehat g$。
\end{theorem}
\begin{proof}
精确 residual 为 $R=U+K-\eta(r+K_{qT})$；替换后 $\widehat R=R-\epsilon_S+\eta\epsilon_T$。因此
\begin{align}
\widehat g-\nabla L_{\mathrm{BPO}}
&=\E_q[(\widehat R-R)\psi]/\eta\\
&=(-\epsilon_S/\eta+\epsilon_T)\E_q[\psi].
\end{align}
由于 $q$ 冻结，
\begin{equation}
\nabla K=-\sum_aq_a\nabla\log p_a=-\E_q[\psi],
\end{equation}
得到第一式。$\epsilon_T$ 冻结，$\nabla L_{\mathrm{BPO}}-\epsilon_T\nabla K$ 的 Jacobian 对称。对剩余乘积逐项求导：
\begin{equation}
\partial_j(\epsilon_S\partial_iK)
=(\partial_j\epsilon_S)(\partial_iK)+\epsilon_S\partial_j\partial_iK.
\end{equation}
交换 $i,j$ 时 Hessian 项抵消，得到所述反对称矩阵。标量势函数的二阶混合偏导相同，所以该矩阵非零排除了固定势函数。
\end{proof}

\begin{example}[三 token 的非保守更新]
\label{opd:curl-example}
取
\begin{equation}
q=(0.5,0.3,0.2),\quad t=(0.2,0.3,0.5),\quad H=\{0\},\quad
\eta=0.7,\quad z=(-0.3,0.6,0),
\end{equation}
并以 $(z_0,z_1)$ 为自由坐标。记 $Q=q_1+q_2=0.5$、$P=p_1+p_2$。粗粒化 student center 为
\begin{equation}
\widehat K_S=q_0\log(q_0/p_0)+Q\log(Q/P).
\end{equation}
完整与粗粒化 KL 之差为
\begin{equation}
\epsilon_S=\sum_{a=1}^2q_a\log(q_a/p_a)-Q\log(Q/P).
\end{equation}
沿 $z_0$ 改变只缩放尾部总质量，尾内比例不变，故 $\partial_0\epsilon_S=0$。逐项微分尾内比例得到
\begin{equation}
\partial_1\epsilon_S=Q\frac{p_1}{P}-q_1,\qquad
\partial_0K=p_0-q_0.
\end{equation}
于是
\begin{equation}
\partial_1\widehat g_0-\partial_0\widehat g_1
=\frac{(p_0-q_0)(Qp_1/P-q_1)}{\eta}
=-0.0095250968566.
\end{equation}
从 log-softmax 计算上述条件更新，其 Jacobian 数值为
\begin{equation}
D\widehat g=
\begin{pmatrix}
0.3491477001&-0.2186486876\\
-0.2091235908&0.2947641411
\end{pmatrix}.
\end{equation}
交叉偏导不同，因此即使动作期望完全枚举、无采样噪声，该粗粒化 BPO 更新也一般没有固定标量目标。完整头部或尾内条件分布完全匹配时，$\epsilon_S=0$，恢复精确梯度。
\end{example}

\begin{proposition}[近似平方回归与 stopped score 的区别]
\label{opd:approximate-square}
令 $\widehat L=\E_q[\widehat R^2]/(2\eta)$，且直接对近似 student center 保留梯度，则
\begin{equation}
\nabla\widehat L-\widehat g
=\frac{\E_q[\widehat R]}{\eta}\nabla\widehat K_S.
\end{equation}
\end{proposition}
\begin{proof}
冻结 reward 和 teacher center 后，$\nabla\widehat R=\psi+\nabla\widehat K_S$。对平方目标使用链式法则：
\begin{align}
\nabla\widehat L
&=\eta^{-1}\E_q[\widehat R\psi]
+\eta^{-1}\E_q[\widehat R]\nabla\widehat K_S\\
&=\widehat g+\eta^{-1}\E_q[\widehat R]\nabla\widehat K_S.
\end{align}
精确中心下 $\E_q[R]=0$，额外项消失；粗粒化中心下该均值为 $-\epsilon_S+\eta\epsilon_T$，一般非零。
\end{proof}

\section{生产缓存、测度与张量契约}
\label{opd:production-contract}

数学 oracle 位于 \path{src/mirroropd/token_inner.py}，支持 Part、Mean、BPO、MDPO 四种总体损失。生产 kernel 位于 \path{src/mirroropd/training/losses.py}，支持 \path{pmd_mean}、\path{bpo}、\path{mdpo}、\path{opd_onpolicy}、\path{opd_is}、\path{opd_is_no_baseline}；生产白名单没有 Part。四方法理论比较与三种 mirror 生产实现须分别报告。

\path{src/mirroropd/training/cache.py} 先用 anchor 完整 logits 构造冻结头部 ids、头部 log-prob、尾部 log-mass、实际 action log-prob 与有效位置 mask。teacher 必须在同一 tokenizer 和 vocabulary 上评分 anchor 选出的 ids，并保留实际 action 的 teacher log-prob。teacher 自己的 Top-$J$ 不能替代 anchor 请求的 ids。两个 worker 顺序执行，传输大小为 prefix 数乘头部大小，完整 logits 可以在各 worker 本地释放。

头部支持集在整个 inner 冻结。teacher center 由头部精确项和尾部质量的单事件 KL 构成，并在 outer 内冻结；student center 在每次 inner forward 重新计算，然后作为 score 权重停止梯度。cache Top-$J$ 表示词表的头尾分割，rollout 的动作支持仍是完整 vocabulary；它不等于截断采样，也不等于限制优化变量的支持。

logits 必须先除以同一 rollout 温度再计算 log-softmax。完整支持语义要求 \path{do_sample=True}、\path{top_p=1}、\path{top_k=-1}，温度严格正且有限。事后重算未截断 log-prob 不能恢复已被截断 proposal 排除的 action；若从行为分布 $\beta$ 采样，目标 $q$ 的回归需要 $q/\beta$ 且 $q\ll\beta$，student 测度目标需要 $p/\beta$ 且 $p\ll\beta$。

verl 适配位于 \path{src/mirroropd/integrations/verl/fsdp.py}。缓存与 logits 按 prediction position 对齐；prompt、padding 和 terminal 无效位置须在 gather 前排除，不能用无效 token id 或 NaN cache 参与计算。适配层负责 padded/jagged 布局及 sequence-parallel 尺寸，损失核心操作本地完整词表而非 tensor-parallel 的独立词表分片。框架、CUDA 和 Ray 耦合留在适配层，生产数学 kernel 可在 CPU 独立导入。

冻结 rollout 在 inner 多步复用。importance ratio 每步依据当前 student 重算，再停止梯度；冻结的是行为策略和样本，不是整轮固定 student ratio。对于参数由同一数据多次更新的过程，总体无偏的单步估计恒等式不能推出“条件于自适应参数和冻结数据仍无偏”。hard sync 在 outer 结束后用训练完成的 student 替换 anchor，再刷新 rollout 与 teacher cache。研究代码的 soft sync 是逐 prefix 概率混合；参数 EMA 是另一映射，生产部署情况需要独立核验。

生产输出的 \path{distillation_losses} 是 stopped coefficient 乘 action log-prob，其数值可正可负，不能当作平方回归或完整 teacher KL 的目标数值。诊断至少区分精确 KL、粗粒化中心、sampled residual、importance ratio、头尾质量、support overlap 和 student entropy。Mirror score 不做静默 ratio clipping，非有限系数会显式报错；外部 gradient clipping 与 ratio clipping具有不同数学语义。

\path{tests/test_upstream_opd_semantics.py} 枚举完整动作核对 full-support 梯度端点、partial-support 中心误差和近似平方目标差异，并验证无效位置的梯度为零；\path{tests/test_upstream_mirror_wiring.py} 核对 anchor-to-teacher 路由、temperature、prediction-position、rollout 配置与版本契约。\path{tests/test_mdpo_eta_limits.py} 核对有限 $\eta$ 梯度、$\eta=1$ teacher KL 和无限 $\eta$ frozen reward 端点。CPU 接口检查、真实模型 canary 与完整训练 benchmark 是不同证据层级。
